% !TeX program = xelatex
% 本地资料蒸馏的工作模板；不是已核实的官方竞赛格式。
% 所有数字须追溯到 result-ledger.csv；按真实题数增删章节。
\documentclass[UTF8,a4paper,zihao=-4,fontset=fandol]{ctexart}
\usepackage[margin=25mm]{geometry}
\usepackage{amsmath,amssymb,bm,booktabs,tabularx,graphicx}
\usepackage{caption,enumitem}
\usepackage[hidelinks]{hyperref}
\IfFontExistsTF{Times New Roman}{\setmainfont{Times New Roman}}{\setmainfont{TeX Gyre Termes}}
\setlength{\parindent}{2em}
\setlength{\parskip}{0pt}
\renewcommand{\baselinestretch}{1.38}
\renewcommand{\arraystretch}{1.2}
\ctexset{
  section={format=\centering\songti\bfseries\zihao{4},beforeskip=1.2ex,afterskip=0.6ex},
  subsection={format=\songti\bfseries\zihao{-4},beforeskip=1ex,afterskip=0.4ex},
  subsubsection={format=\songti\bfseries\zihao{-4},beforeskip=0.8ex,afterskip=0.4ex}
}
\captionsetup{font=small,labelfont=bf}
\setlist{nosep,leftmargin=2em}
\newcolumntype{Y}{>{\raggedright\arraybackslash}X}
\newcommand{\slot}[1]{〔待填：#1〕}
\hypersetup{pdfauthor={},pdftitle={数学建模论文工作模板}}
\pagestyle{plain}
\raggedbottom
\clubpenalty=10000
\widowpenalty=10000

\begin{document}
\begin{center}
{\songti\bfseries\fontsize{16pt}{20pt}\selectfont \slot{研究对象与核心问题}}\par
\vspace{1em}
{\heiti\bfseries\zihao{4} 摘要}
\end{center}

\slot{用一至两句界定研究对象、任务与关键约束。}

\textbf{针对问题\slot{题号}，}\slot{根据什么数据、抓住什么关系、建立什么模型、采用什么求解方法，得到什么答案。填写单位、范围及真实验证结果；草稿中注明结果台账 ID。}

% 按实际问题数复制上一段；不要保留不存在的题号。
\slot{总结最有证据支持的发现和适用边界。没有执行的检验不能写成已完成。}

\noindent\textbf{关键词：}\slot{研究对象}；\slot{核心模型}；\slot{求解方法}

% 摘要另页沿用本地模板习惯，可按适用竞赛规则修改。
\clearpage
\section{问题重述}
\slot{题目提供的对象、输入、时空范围和关键约束；区分事实与假设。}
\subsection{问题与交付物}
\slot{逐问列明需要输出的数值、方案、预测或解释，保留单位与评价口径。}

\section{问题分析}
\slot{说明各问目标、难点、方法选择理由及输入输出依赖；流程图仅在能说明实际计算过程时使用。}

\section{模型假设}
\slot{依据假设台账说明每项简化、理由、影响和验证方法。假设数量由问题决定。}

\section{符号说明}
\begin{table}[htbp]
\centering
\caption{主要符号及单位}
\label{tab:symbols}
\small
\begin{tabularx}{\linewidth}{lYY}
\toprule
符号 & 含义 & 单位或取值域 \\
\midrule
\slot{符号} & \slot{实际变量含义} & \slot{单位或集合} \\
\bottomrule
\end{tabularx}
\end{table}

\section{问题\slot{题号}的模型建立与求解}
\subsection{数据与建模依据}
\slot{输入文件和字段、单位、处理规则、处理前后样本量、模型选择依据。}

\subsection{模型建立}
\slot{由现实关系推导公式，定义变量、目标、约束、参数来源和边界条件。公式与代码口径一致。}
% 公式用 equation/align 环境，按问题设置唯一标签，如 eq:q1-model；
% 正文以 \eqref{eq:q1-model} 引用，复制章节时相应改为 eq:q2-model 等。
% 此处不预置具体模型，避免将模板形式当作本题模型。

\subsection{求解方法与实现}
\slot{算法、关键参数、初始条件、停止准则、随机种子、运行环境和脚本入口。准确报告求解状态。}

\subsection{结果与分析}
\slot{先给本问关键答案，再解释图表。每个数字、单位与口径均对应真实结果台账 ID。}
\begin{table}[htbp]
\centering
\caption{本问核心结果}
\label{tab:q1-results}
\small
\begin{tabularx}{\linewidth}{YYYY}
\toprule
输出 & 结果及单位 & 条件与解释 & 证据 ID \\
\midrule
\slot{输出项} & \slot{真实结果} & \slot{条件与题意} & \slot{R-ID} \\
\bottomrule
\end{tabularx}
\end{table}

\subsection{模型验证}
\slot{选择与模型风险相匹配的验证：样本外误差、约束可行性、量纲守恒、基线比较、权重稳定性或随机误差。填写实际方法、范围、结果和证据 ID。}

% 按实际问题数复制完整模型章节，LaTeX 自动连续编号。
% 同时把 tab:q1-results 等标签改成 tab:q2-results，并更新对应引用；
% 公式、图、表、章节的 label 都必须唯一，不能原样复制同名标签。
\section{模型评价与改进}
\textbf{有证据的优点。}\slot{设计选择带来什么收益，由什么结果支持。}

\textbf{主要局限。}\slot{未覆盖条件、潜在偏差及受影响结论。}

\textbf{改进路径。}\slot{针对局限需要新增的数据、模型或检验；区分已经完成与未来建议。}

\section{结论}
\slot{逐问回收答案及适用条件，不引入正文没有的新数字或发现。}

\section*{参考文献}
\slot{仅列实际查阅、核验并在正文引用的来源。}
% 启用 BibTeX 时同时移除上方手写的 \section*{参考文献} 及占位段落，
% 避免自动 bibliography 再次生成参考文献标题；再按适用格式取消以下注释。
% \bibliographystyle{gbt7714-numerical}
% \bibliography{references}
% 只使用实际 \cite{key}，不使用 \nocite{*}。

\appendix
\section{支撑文件与复现入口}
\slot{真实文件列表、功能、输入输出、运行顺序、依赖版本，以及确实验证过的命令。}

\section{AI 使用记录}
\slot{按真实过程记录工具和版本、用途、关键交互、采纳修改及核验证据。是否需要本节及其位置以适用官方通知为准。}

% 提交前清除所有填写说明与占位符，逐页检查排版。
\end{document}
